% Part: incompleteness
% Chapter: theories-computability
% Section: computably-axiomatizable

\documentclass[../../../include/open-logic-section]{subfiles}

\begin{document}

\olfileid[tr]{inc}{tcp}{cax}

\olsection{\printtoken{S}{axiomatizable} Kuramlar}

Bir $\Th{T}$ kuramına, hesaplanabilir bir $A$ aksiyomlar kümesine sahipse
\emph{!!{axiomatizable}} denir. ($A$'nın $\Th{T}$ için bir aksiyomlar kümesi
olduğunu söylemek, $T = \Setabs{!A}{A \Proves !A}$ demektir.) Doğal
sayıların her “makul” aksiyomlaştırması bu özelliğe sahip olacaktır.
Özellikle sonlu bir aksiyomlar kümesine sahip her kuram
!!{axiomatizable}{}dir.

\begin{lem}
$\Th{T}$'nin !!{axiomatizable} olduğunu varsayın. O zaman $\Th{T}$,
!!{computably
enumerable}{}dir.
\end{lem}

\begin{proof}
$A$, $\Th{T}$ için hesaplanabilir bir aksiyomlar kümesi olsun. $!A \in T$
olup olmadığını belirlemek için aksiyomlardan $!A$ için !!a{derivation}
aramak yeterlidir.

Biraz başka türlü söylersek $!A$, $\Th{T}$ içindedir ancak ve ancak $A$
içinde $!B_1$, \dots, $!B_k$ biçiminde sonlu bir aksiyomlar listesi ve
birinci derece mantıkta
$(!B_1 \land \dots \land !B_k) \lif !A$ için !!a{derivation} vardır.
Fakat ikinci “\dots” kısmı hesaplanabilir olmak koşuluyla, “öyle bir \dots
vardır ki \dots” biçiminde tanımı bulunan her kümenin !!{c.e.} olduğunu
zaten biliyoruz.
\end{proof}

\end{document}
